% Vietnamese translation for OI Public Library
% Source-File: IOI中国国家候选队论文/2004/何林.ppt
\documentclass[11pt]{article}
\usepackage{oipl-vn}

\newcommand{\Oh}{\mathcal{O}}

\title{Phương pháp bảo toàn trong tin học\\{\large Ghi chú theo slide}}
\author{He Lin}
\date{IOI 2004}

\begin{document}
\maketitle

\origtitle{Conservation method in informatics}
\sourcepath{IOI 2004 / He Lin.ppt}

\section{Đường dây của bài trình bày}

Tệp trình chiếu là PPT nhị phân cũ nên phần văn bản đọc được chủ yếu là các
mốc công thức và ví dụ. Ghi chú này bám vào các mốc ấy, đồng thời dựa trên
bản dịch bài viết PDF cùng chủ đề.

Chủ đề trung tâm là \term{bảo toàn}: khi một quá trình biến đổi có quá nhiều
chi tiết, ta tìm đại lượng không đổi để đi thẳng tới kết luận. Hai ví dụ mở
đầu trên slide là va chạm đàn hồi với vận tốc \(5\,\mathrm{m/s}\),
\(4\,\mathrm{m/s}\), và phản ứng giữa \(10\mathrm{g}\) carbon với
\(10\mathrm{g}\) \(O_2\) trong bình kín. Chúng nhắc rằng động lượng, động
năng hoặc khối lượng có thể quan trọng hơn mô phỏng từng bước.

\section{Ví dụ dãy số}

Một mốc slide nổi bật là phép biến đổi cục bộ trên ba phần tử liên tiếp:
\[
  (a_{i-1},a_i,a_{i+1})
  \longmapsto
  (a_{i-1}+a_i,-a_i,a_i+a_{i+1}).
\]
Nếu đặt tổng tiền tố \(S_i=a_1+\cdots+a_i\), phép biến đổi này hoán đổi hai
tổng tiền tố lân cận \(S_{i-1}\) và \(S_i\). Vì vậy đa tập
\(\{S_1,S_2,\ldots,S_{n-1}\}\) được bảo toàn theo nghĩa chỉ thay đổi thứ tự.
Từ bất biến đó, ta có thể xét khả năng biến đổi giữa hai dãy bằng cách so
sánh các tổng tiền tố, thay vì lần theo mọi chuỗi thao tác.

Slide cũng có ví dụ dữ liệu vào ra và nhắc độ phức tạp \(\Oh(n\log n)\), phù
hợp với cách sắp xếp rồi so sánh các giá trị tổng tiền tố.

\section{Thông điệp}

Phương pháp bảo toàn đặc biệt hữu ích khi thao tác địa phương gây ra hiệu ứng
toàn cục khó quan sát. Một khi tìm được bất biến đúng, bài toán thường chuyển
từ mô phỏng sang kiểm chứng điều kiện cần và đủ. Bài trình bày kết thúc bằng
một ví dụ liên quan đến Fibonacci, tiếp tục nhấn mạnh cùng nguyên tắc: chọn
đại lượng được giữ nguyên trước khi chọn thuật toán.

\end{document}
